% Study sheet: Volume IV, Chapter 1, The Causal Boundary.
% Numbers: \val{...} from boundary_studysheet_numbers.tex, written by boundary_studysheet_numbers.py.
\documentclass{studysheet}
\input{boundary_studysheet_numbers}
\sheetvolume{4}
\sheetchapter{1}{The Causal Boundary}
\sheetsources{Numbers from the chapter's running machine (illustrative) and the NTSB Tempe record.}

\begin{document}
% ===================================================================== FRONT: LEARN
\frontside
\begin{bigidea}
A learned model's output is only a \hl{proposal} until it crosses the \hl{causal boundary}, where a command
becomes current, force and motion. Work delivered there cannot be recalled, and the machine keeps moving
while the model computes. An independent \hl{permission path}, not the model, decides what may cross.
\end{bigidea}

\begin{leftcol}
\sheetsec{Key ideas}
\begin{ideas}
\idea{Proposal, then work.}{Before the boundary a tensor is discarded for free; after it, changing a
  command costs work against inertia.}
\idea{Embodiment loses three privileges.}{Unlike advisory or digital ML: no paused time (delay is distance), no rollback (recovery is a
  new motion), no batching.}
\idea{The scope test.}{Learned model, physical feedback, delegated authority. Per-command
  approval removes authority; a veto or a planner after the model does not.}
\idea{Proposal is not permission.}{The Brain proposes; the permission path admits only with fresh state,
  limits met and a fallback ready, and runs that fallback.}
\idea{Independence includes evidence.}{A check fed by the proposer's cameras fails with them; cross-check plant
  signals: wheel slip, IMU acceleration, motor current.}
\idea{Three demands meet there.}{Fresh, authenticated state; authority to veto and run a fallback;
  a synchronized record of proposal, decision, command and response.}
\end{ideas}

\vspace{10pt}
\sheetsec{The four bedrock laws}
\begin{sheettable}{@{}L l L@{}}
\thead{Law} & \thead{Part} & \thead{In the field: what it demands}\\ \midrule
First: irreversibility & I anatomy & irreversible work: check before, not after\\
Second: endogenous data & II teaching & covariate shift: offline scores do not certify\\
Third: proposal is not permission & III running & runtime assurance (Simplex): an independent check each tick\\
Fourth: permission must be earned & IV governing & assurance case, ODD: authority only where evidence covers\\
\end{sheettable}
\vspace{3pt}
{\fontsize{8.6}{11}\selectfont\color{muted}Only the first follows from physics; the book makes the other
three mandatory by design.\par}

\end{leftcol}%
\begin{rightcol}
\sheetsec{The five-level machine}
{\centering
\begin{tikzpicture}[sheet]
  \node[accbox, minimum width=150pt, minimum height=15pt] (br) at (75,112) {Brain: learned proposers, \val{intenthz}--\val{stephz}\,Hz};
  \draw[accent, dashed, line width=0.7pt] (-6,98) -- (156,98);
  \node[nbox, minimum width=150pt, minimum height=15pt] (ns) at (75,84) {Nervous System: permission path, \val{tickkhz}\,kHz};
  \node[nbox, minimum width=150pt, minimum height=15pt] (bo) at (75,62) {Body: drives, windings, linkages};
  \node[hotbox, minimum width=150pt, minimum height=15pt] (bd) at (75,40) {Boundary: register write, sensor capture};
  \node[nbox, dashed, minimum width=150pt, minimum height=15pt] (w) at (75,18) {World: floor, payload, people};
  \node[note, anchor=west, align=left] at (162,98) {proposal boundary:\\who may approve};
  \node[note, anchor=west, align=left] at (162,40) {causal boundary:\\where work begins};
\end{tikzpicture}\par}
\vspace{4pt}
{\fontsize{8.6}{11}\selectfont Levels run from computation down to physics, not along the command path.
Governance wraps all five; it is not a sixth level.\par}

\vspace{4pt}
\fcard{Elementary stopping budget}%
  {\StoppingDistance = \SpeedAtDetection\PrebrakeDelay + \frac{\SpeedAtDetection^2}{2\CredibleBrakingDeceleration} \le \ClearDistance,
   \qquad \StoppingCrossoverSpeed = 2\CredibleBrakingDeceleration\PrebrakeDelay}%
  {Speed $\SpeedAtDetection$ times pre-brake delay $\PrebrakeDelay$, plus braking at $\CredibleBrakingDeceleration$,
   must fit the clear distance $\ClearDistance$. Below crossover speed $\StoppingCrossoverSpeed$ delay travel is larger.}%
  {Halving speed halves delay travel and quarters braking. Drive limit, to scale:}%
  {\begin{tikzpicture}[sheet]
     \pgfmathsetmacro\sc{226/\val{clearmm}}
     \pgfmathsetmacro\xa{\val{delaytravel}*\sc}
     \pgfmathsetmacro\xb{(\val{delaytravel}+\val{braketravel})*\sc}
     \draw[accbox] (0,0) rectangle (\xa,12);
     \draw[hotbox] (\xa,0) rectangle (\xb,12);
     \draw[dim, dashed] (\xb,0) rectangle (226,12);
     \draw[accent, line width=1pt] (226,-3) -- (226,15);
     \node[small, anchor=north] at ({\xa/2},-3) {delay \val{delaytravel}};
     \node[small, anchor=north] at ({(\xa+\xb)/2},-3) {braking \val{braketravel}};
     \node[small, anchor=north] at ({(\xb+226)/2},-3) {spare \val{spare} mm};
     \node[lab, anchor=south east] at (226,15) {$\ClearDistance$ \val{clearmmc}\,mm};
   \end{tikzpicture}}

\vspace{6pt}
{\color{rulegray}\hrule height 0.5pt}\vspace{6pt}
{\sheetsans\fontsize{7.4}{9}\selectfont\bfseries\color{muted}PRE-BRAKE DELAY: FIVE TERMS, TO SCALE (MS)}\par\vspace{5pt}
{\centering
\begin{tikzpicture}[sheet]
  \pgfmathsetmacro\sc{226/\val{tau}}
  \pgfmathsetmacro\xb{\val{backbone}*\sc}
  \pgfmathsetmacro\xa{\val{age}*\sc}
  \pgfmathsetmacro\xl{(\val{age}+\val{lease})*\sc}
  \pgfmathsetmacro\xt{(\val{age}+\val{lease}+\val{tick}+\val{bus})*\sc}
  \draw[hotbox] (0,0) rectangle (\xb,12);
  \draw[accbox] (\xb,0) rectangle (\xa,12);
  \draw[nbox] (\xa,0) rectangle (\xl,12);
  \draw[accbox] (\xl,0) rectangle (\xt,12);
  \draw[nbox] (\xt,0) rectangle (226,12);
  \node[small, anchor=south west] at (0,15) {learned backbone \val{backbone}, inside the age};
  \draw[dim] ({(\xl+\xt)/2},12.5) -- ({(\xl+\xt)/2},14.5);
  \node[small, anchor=south] at ({(\xl+\xt)/2},19.5) {tick + bus \val{tick} + \val{bus}};
  \draw[dim] (0.5,-2) -- (0.5,-5) -- ({\xa-0.5},-5) -- ({\xa-0.5},-2);
  \node[small, anchor=north] at ({\xa/2},-7) {observation age \val{age}};
  \node[small, anchor=north] at ({(\xa+\xl)/2},-7) {chunk lease \val{lease}};
  \node[small, anchor=north] at ({(\xt+226)/2},-7) {onset \val{onset}};
\end{tikzpicture}\par}
\vspace{9pt}
{\fontsize{8.6}{11}\selectfont Sum: $\PrebrakeDelay$ = \val{tau}\,ms. A stalled proposer's chunk drives on until its lease expires.
Each millisecond costs \val{permm}\,mm at \val{speed}\,m/s.\par}

\vspace{8pt}
\sheetsec{Three machine classes: the margin each spends}
\begin{sheettable}{@{}l L@{}}
\thead{Class} & \thead{Spends}\\ \midrule
Mass-dominated mobility & stopping clearance (momentum)\\
Contact-dominated manipulation & force and deformation margin\\
Energy-, thermal-dominated & heat and supply margin\\
\end{sheettable}
\end{rightcol}

% ===================================================================== BACK: PRACTICE
\backside
\begin{answers}
\ans{Finish}{$\val{speed} \times \val{taubigs}$ = \val{delaybig}\,mm; stop \val{stopbig}\,mm,
  \val{sparebig}\,mm spare: it fits, \val{addedmm}\,mm costlier. $\StoppingCrossoverSpeed$ = \val{twodecel}
  $\times$ \val{taubigs} $\approx$ \val{crossoverbig}\,m/s, below \val{speed}: braking dominates.}
\ans{1}{(a) Mean \val{tempemean}\,m/s $\times$ \val{tempesec}\,s $\approx$ \res{\val{tempedist}\,m}: delay
  travel, never recovered (first law). (b) Dropped history: fresh state. Suppression and braking off: authority
  to intervene. Unwatched operator: authority on evidence nobody collected.}
\columnbreak
\ans{2}{Highlighter: a person approves each removal. Governor: no learned model. Base: physical AI; a veto
  does not remove delegated authority (scope test). Travel: stopping clearance; latch: force margin; hold:
  heat and supply margin.}
\ans{3}{(a) \val{stall}\,ms $\times$ \val{permm}\,mm/ms = \res{\val{stallmm}\,mm}. (b) Lease expiry at \val{lease}\,ms
  (\val{leasemm}\,mm): the permission path (Nervous System) starts the fallback stop, then onset and braking.
  Shared cameras fail together.}
\ans{4}{(a) No: \val{muoil} is below \val{muneed}. (b) \val{muoil} $\times$ \val{gravity} $\approx$
  \val{tractionceiling}\,m/s²; \val{speedsq} $\div$ \val{twoceiling} $\approx$ \res{\val{brakeslick}\,m},
  far past \val{clear}\,m. (c) Wheel slip, commanded
  versus IMU acceleration, current transients.}
\end{answers}

\begin{leftcol}
\sheetsec{Worked example}
\begin{practicebox}\boxtitle{Does the loaded base stop inside its clearance?}
Toward a rack end at the drive limit, $\SpeedAtDetection$ = \val{speed}\,m/s, with
$\CredibleBrakingDeceleration$ = \val{decel}\,m/s², $\ClearDistance$ = \val{clear}\,m, $\PrebrakeDelay$ =
\val{tau}\,ms. Which term dominates?
\begin{steps}
\item Delay travel: $\val{speed}\,\text{m/s} \times \val{taus}\,\text{s}$ = \res{\val{delaytravel}\,mm}.
\item Braking: $\val{speedsq} \div (\val{twodecel}\,\text{m/s}^2)$ = \res{\val{braketravel}\,mm}.
\item Stop: \val{delaytravel} + \val{braketravel} = \res{\val{stopmm}\,mm} $\le$ \val{clearmm}\,mm,
  with \val{spare}\,mm spare.
\item Crossover: $\val{twodecel} \times \val{taus}$ $\approx$ \res{\val{crossover}\,m/s}; the drive limit is
  above it, so braking dominates (\val{brakeratio}$\times$ the delay travel).
\end{steps}
\takeaway{The backbone costs \val{backbonemm}\,mm, the lease \val{leasemm}\,mm. The spare is not yet a
safety margin; later allowances must fit in it.}
\end{practicebox}

\vspace{10pt}
\sheetsec{Your turn: finish it}
\begin{practicebox}\boxtitle{How much clearance does a slower, more accurate detector cost?}
\assumed{} a more accurate backbone takes \val{bigbackbone}\,ms instead of
\val{backbone}\,ms.\par\vspace{4pt}
\textbf{Given:} $\PrebrakeDelay$ = \val{tau} $-$ \val{backbone} + \val{bigbackbone} = \val{taubig}\,ms.\par
\textbf{You:} find the new stop and spare, and whether braking still dominates.\par\vspace{56pt}
\end{practicebox}

\end{leftcol}%
\begin{rightcol}
\sheetsec{Check yourself}
\begin{qs}
\item Tempe: each change of the pedestrian's label dropped her earlier positions. Recognizing an emergency at
  \val{tempea}\,m/s, the vehicle withheld braking for \val{tempesec}\,s, slowing to \val{tempeb}\,m/s. Factory
  emergency braking was off; nobody watched the operator. (a) Distance with braking withheld, at the mean
  speed? (b) Which of the three demands does each failure break?
\item Physical AI or not, and which criterion decides? (i) A learned detector flags faulty parts on a
  screen. (ii) A PID governor holds conveyor speed. (iii) The base: a learned detector feeds a
  hand-written planner; a certified scanner can cut the drives. Which margin does the base spend as it
  travels, presses a latch, holds a load?
\item At the drive limit the chunk policy stalls for \val{stall}\,ms just after a renewal. (a) With no lease,
  how far would the base drive on its last command? (b) What ends the motion, at which level, and when? Why must that check
  not use the proposer's cameras?
\item \assumed{} hydraulic fluid drops floor friction from \val{mudry} to \val{muoil}; the card's deceleration
  assumes at least \val{muneed}, and traction caps deceleration at friction $\times$ \val{gravity}\,m/s².
  (a) Does the card still hold? (b) Least braking travel from the drive limit? (c) The cameras see nothing: what
  reveals the slick?
\end{qs}

\vspace{8pt}
\sheetsec{Common mistakes}
\mistake{``Add a collision penalty to the reward and the policy learns to be safe.''}%
  {A loss improves average training behavior; it guarantees nothing out of distribution. A hard check
   runs every tick, downstream (third law).}{}
\mistake{``If the robot misbehaves, the operator takes over.''}%
  {In-loop takeover takes \val{takeover}\,ms: \val{takeoverticks} permission ticks and
   \val{takeovermm}\,mm at the drive limit, \val{takeovershare}\% of the clear distance.}{}
\end{rightcol}
\end{document}
